\usepackage{amsthm,mdframed}
\usepackage{nameref}

%% FROM https://texblog.org/2015/09/30/fancy-boxes-for-theorem-lemma-and-proof-with-mdframed/
\newcounter{theo}[chapter]
\setcounter{theo}{0}
\renewcommand{\thetheo}{\thechapter.\arabic{theo}}

\labelformat{theo}{\thechapter-#1}


%%%%%%%%%%%%%%%%%% TEOREMA
% Following https://tex.stackexchange.com/questions/705270/mdthereom-shared-counters-cleverref-and-nested-theorems
\let\theorem\relax

\newcounter{theorem}
\makeatletter
\let\c@theorem\c@theo
\let\thetheorem\thetheo
\makeatother

\makeatletter
\newenvironment{theorem}[1][]{%
        % \protected@edef\@currentlabelname{#1}
        \refstepcounter{theorem}%
        \ifstrempty{#1}%
        {\mdfsetup{%
                        frametitle={%
                                        \tikz[baseline=(current bounding box.east),outer sep=0pt]
                                        \node[anchor=east,rectangle,fill=orange!10]
                                        {\strut Theorem~\thetheo};}}
        }%
        {
                \mdfsetup{%
                        frametitle={%
                                        \tikz[baseline=(current bounding box.east),outer sep=0pt]
                                        \node[anchor=east,rectangle,fill=orange!20]
                                        {\strut Theorem~\thetheo:~#1};}}%
        }%
        \mdfsetup{innertopmargin=10pt,linecolor=orange!10,%
                linewidth=2pt,topline=true,%
                frametitleaboveskip=\dimexpr-\ht\strutbox\relax
        }
        \begin{mdframed}[]}
                {\end{mdframed}}
\makeatother
\Crefname{theorem}{Theorem}{Theorems}

%%%%%%%%%%%%%%%%%% TEOREMA
\newcounter{strongtheorem}
\makeatletter
\let\c@strongtheorem\c@theo
\let\thestrongtheorem\thetheo
\makeatother

\makeatletter
\newenvironment{strongtheorem}[1][]{%
        % \protected@edef\@currentlabelname{#1}
        \refstepcounter{strongtheorem}%
        \ifstrempty{#1}%
        {\mdfsetup{%
                        frametitle={%
                                        \tikz[baseline=(current bounding box.east),outer sep=0pt]
                                        \node[anchor=east,rectangle,fill=red!20]
                                        {\strut Theorem~\thetheo};}}
        }%
        {
                \protected@edef\@currentlabelname{#1}
                \mdfsetup{%
                        frametitle={%
                                        \tikz[baseline=(current bounding box.east),outer sep=0pt]
                                        \node[anchor=east,rectangle,fill=red!20]
                                        {\strut Theorem~\thetheo:~#1};}}%
        }%
        \mdfsetup{innertopmargin=10pt,linecolor=red!20,%
                linewidth=2pt,topline=true,%
                frametitleaboveskip=\dimexpr-\ht\strutbox\relax
        }
        \begin{mdframed}[]}
                {\end{mdframed}}
\makeatother
\Crefname{strongtheorem}{Theorem}{Theorems}

%%%%%%%%%%%%%%%%%% LEMA
\newcounter{lemma}
\makeatletter
\let\c@lemma\c@theo
\let\thelemma\thetheo
\makeatother

\makeatletter
\newenvironment{lemma}[1][]{%
        \refstepcounter{lemma}%
        \ifstrempty{#1}%
        {\mdfsetup{%
                        frametitle={%
                                        \tikz[baseline=(current bounding box.east),outer sep=0pt]
                                        \node[anchor=east,rectangle,fill=orange!20]
                                        {\strut Lemma~\thetheo};}}
        }%
        {
                \protected@edef\@currentlabelname{#1}
                \mdfsetup{%
                        frametitle={%
                                        \tikz[baseline=(current bounding box.east),outer sep=0pt]
                                        \node[anchor=east,rectangle,fill=orange!20]
                                        {\strut Lemma~\thetheo:~#1};}}%
        }%
        \mdfsetup{innertopmargin=10pt,linecolor=orange!20,%
                linewidth=2pt,topline=true,%
                frametitleaboveskip=\dimexpr-\ht\strutbox\relax
        }
        \begin{mdframed}[]}
                {\end{mdframed}}
\makeatother
\Crefname{lemma}{Lemma}{Lemmas}

%%%%%%%%%%%%%%%%%% DEFINICION
\newcounter{definition}
\makeatletter
\let\c@definition\c@theo
\let\thedefinition\thetheo
\makeatother

\makeatletter
\newenvironment{definition}[1][]{%
        \refstepcounter{definition}%
        \ifstrempty{#1}%
        {\mdfsetup{%
                        frametitle={%
                                        \tikz[baseline=(current bounding box.east),outer sep=0pt]
                                        \node[anchor=east,rectangle,fill=blue!20]
                                        {\strut Definition~\thetheo};}}
        }%
        {
                \protected@edef\@currentlabelname{#1}
                \mdfsetup{%
                        frametitle={%
                                        \tikz[baseline=(current bounding box.east),outer sep=0pt]
                                        \node[anchor=east,rectangle,fill=blue!20]
                                        {\strut Definition~\thetheo:~#1};}}%
        }%
        \mdfsetup{innertopmargin=10pt,linecolor=blue!20,%
                linewidth=2pt,topline=true,%
                frametitleaboveskip=\dimexpr-\ht\strutbox\relax
        }
        \begin{mdframed}[]}
                {\end{mdframed}}
\makeatother

\Crefname{definition}{Definition}{Definitions}


%%%%%%%%%%%%%%%%%% REMARK
\newcounter{remark}
\makeatletter
\let\c@remark\c@theo
\let\theremark\thetheo
\makeatother

\makeatletter
\newenvironment{remark}[1][]{%
        \refstepcounter{remark}%
        \ifstrempty{#1}%
        {\mdfsetup{%
                        frametitle={%
                                        \tikz[baseline=(current bounding box.east),outer sep=0pt]
                                        \node[anchor=east,rectangle,fill=gray!20]
                                        {\strut Remark~\thetheo};}}
        }%
        {
                \protected@edef\@currentlabelname{#1}
                \mdfsetup{%
                        frametitle={%
                                        \tikz[baseline=(current bounding box.east),outer sep=0pt]
                                        \node[anchor=east,rectangle,fill=gray!20]
                                        {\strut Remark~\thetheo:~#1};}}%
        }%
        \mdfsetup{innertopmargin=10pt,linecolor=gray!20,%
                linewidth=2pt,topline=true,%
                frametitleaboveskip=\dimexpr-\ht\strutbox\relax
        }
        \begin{mdframed}[]}
                {\end{mdframed}}
\makeatother

\Crefname{remark}{Remark}{Remarks}

%%%%%%%%%%%%%%%%%% EXAMPLE
\newcounter{example}
\makeatletter
\let\c@example\c@theo
\let\theexample\thetheo
\makeatother

\makeatletter
\newenvironment{example}[1][]{%
        \refstepcounter{example}%
        \ifstrempty{#1}%
        {\mdfsetup{%
                        frametitle={%
                                        \tikz[baseline=(current bounding box.east),outer sep=0pt]
                                        \node[anchor=east,rectangle,fill=green!20]
                                        {\strut Example~\thetheo};}}
        }%
        {
                \protected@edef\@currentlabelname{#1}
                \mdfsetup{%
                        frametitle={%
                                        \tikz[baseline=(current bounding box.east),outer sep=0pt]
                                        \node[anchor=east,rectangle,fill=green!20]
                                        {\strut Example~\thetheo:~#1};}}%
        }%
        \mdfsetup{innertopmargin=10pt,linecolor=green!20,%
                linewidth=2pt,topline=true,%
                frametitleaboveskip=\dimexpr-\ht\strutbox\relax
        }
        \begin{mdframed}[]}
                {\end{mdframed}}
\makeatother

% \mdfsetup{skipabove=1em,skipbelow=0em}
\surroundwithmdframed[
        topline=false,
        rightline=false,
        bottomline=false,
        leftmargin=0em,%\parindent,
        linewidth=4pt,
        linecolor=gray!40,
        skipabove=\medskipamount,
        skipbelow=\medskipamount
]{proof}


\newcommand{\namedCref}[1]{\Cref{#1}~(\nameref{#1})}